\documentclass[12pt,a4paper]{article}

\usepackage[utf8]{inputenc}
\usepackage[T1]{fontenc}
\usepackage{lmodern}
\usepackage[margin=2.5cm]{geometry}
\usepackage{amsmath,amssymb}
\usepackage{booktabs}
\usepackage{hyperref}
\usepackage{listings}
\usepackage{xcolor}
\usepackage{graphicx}
\usepackage{caption}
\usepackage{parskip}
\usepackage{microtype}
\usepackage{titlesec}
\usepackage{enumitem}

% ── code listing style ──────────────────────────────────────────────────────
\definecolor{codebg}{HTML}{F6F8FA}
\definecolor{codecomment}{HTML}{6A737D}
\definecolor{codestring}{HTML}{032F62}
\definecolor{codekw}{HTML}{D73A49}

\lstdefinestyle{pycode}{
  language=Python,
  backgroundcolor=\color{codebg},
  basicstyle=\ttfamily\small,
  keywordstyle=\color{codekw}\bfseries,
  commentstyle=\color{codecomment}\itshape,
  stringstyle=\color{codestring},
  breaklines=true,
  frame=single,
  framerule=0.4pt,
  rulecolor=\color{gray!40},
  xleftmargin=1em,
  xrightmargin=1em,
  captionpos=b,
  numbers=none,
}

\lstdefinestyle{tscode}{
  language=Java,
  backgroundcolor=\color{codebg},
  basicstyle=\ttfamily\small,
  keywordstyle=\color{codekw}\bfseries,
  commentstyle=\color{codecomment}\itshape,
  stringstyle=\color{codestring},
  breaklines=true,
  frame=single,
  framerule=0.4pt,
  rulecolor=\color{gray!40},
  xleftmargin=1em,
  xrightmargin=1em,
  captionpos=b,
  numbers=none,
}

% ── section spacing ──────────────────────────────────────────────────────────
\titlespacing*{\section}{0pt}{1.4ex plus .2ex}{0.8ex plus .1ex}
\titlespacing*{\subsection}{0pt}{1.2ex plus .2ex}{0.6ex plus .1ex}
\titlespacing*{\subsubsection}{0pt}{1.0ex plus .2ex}{0.4ex plus .1ex}

% ── hyperref setup ───────────────────────────────────────────────────────────
\hypersetup{
  colorlinks=true,
  linkcolor=blue!60!black,
  urlcolor=blue!70!black,
  citecolor=green!50!black,
}

% ─────────────────────────────────────────────────────────────────────────────
\title{\textbf{EEG Based Trigger Prediction and\\Real Time BCI System}\\[0.5em]
       \large Technical Report}
\author{Vinicius Lima}
\date{May 2026}
% ─────────────────────────────────────────────────────────────────────────────

\begin{document}

\maketitle
\tableofcontents
\newpage

% =============================================================================
\section{Objective}
% =============================================================================

This project aims to build a real time Brain Computer Interface (BCI)
system that predicts, from multi channel EEG signals alone, whether a
subject is \emph{about to press a physical trigger button}. The ultimate
application is motor intention decoding: the classifier must issue a
prediction at the \emph{next} sample given only a finite causal look back
window of raw EEG, so that the output can drive downstream actions without
any future data leakage.

As a controlled development and validation step, the pipeline is trained
and evaluated on an eye open/closed detection task - a well characterised
paradigm that generates reliable, high SNR EEG correlates and provides a
ground truth label at every sample. Demonstrating high accuracy on this
proxy task validates the causal windowing scheme, the feature
representation, and the online inference architecture before transitioning
to button press data.

The specific objectives are:

\begin{enumerate}[leftmargin=1.5em]
  \item Train a causal classifier on a recorded EEG dataset (eye open/closed
        as test proxy) using a sliding window feature representation that
        strictly respects temporal causality.
  \item Tune hyperparameters automatically, including the optimal look back
        window length~$T$.
  \item Integrate the trained model into the PiEEG server acquisition
        pipeline for sample by sample online inference.
  \item Extend the server hardware layer to capture a physical button
        signal alongside EEG, providing the ground truth label for
        trigger press prediction during live sessions.
  \item Stream the real time prediction probability time series to the web
        dashboard for monitoring.
\end{enumerate}

% =============================================================================
\section{Methods}
% =============================================================================

\subsection{Dataset}

Recording sessions were collected with a 16 channel PiEEG headset at a
sampling rate of 250\,Hz \textbf{(128 Hz for the eyes open/closed dataset; I might have used this value to convert window length to time so it could have errors bellow)}. As a proxy task for validating the causal
prediction pipeline, the \texttt{eyeDetection} binary label is used during
training: 0~=~eyes closed, 1~=~eyes open. This label serves as a
convenient stand in for button press prediction because eye open/closed
transitions produce strong, well understood EEG correlates, making classification
performance a reliable indicator of whether the causal windowing and
feature extraction scheme is working correctly.

The dataset used for training (\texttt{session\_test\_2.csv}) contains
21\,318 samples with a class distribution of 13\,003 closed and 8\,315
open, yielding a moderate class imbalance of approximately 61\% / 39\%.
Each CSV file contains 16 EEG channel columns (\texttt{ch1}–\texttt{ch16})
and a \texttt{timestamp} column.

\subsection{Preprocessing}

Two preprocessing steps are applied before feature extraction.

\subsubsection{Artifact Removal}

Samples exceeding the 0.1st or 99.9th quantile per channel are flagged as
outliers. A margin of $\pm5$ samples around each flagged point is also
marked to capture the rising and falling edges of transient artifacts. All
flagged values are replaced by linear interpolation along the time axis.
This step replaced 1\,243 sample channel values in the training session.

\subsubsection{Band Pass Filtering}

A zero phase FIR band pass filter (0.1–60\,Hz, Hamming window) is applied
using MNE's \texttt{filter\_data} to remove DC drift and high frequency
noise above the Nyquist relevant band.

\subsubsection{Standardisation}

After filtering, each channel is z score normalised to zero mean and unit
variance using scikit learn's \texttt{StandardScaler}. The fitted scaler is
serialised to disk for use during online inference.

\subsection{Causal Random Forest Classifier}

\subsubsection{Motivation and Causal Window Design}

The model must operate in a strictly causal setting: at inference time only
past and present samples are available. The classification task is:

\begin{quote}
  Given $T$ consecutive EEG samples ending at time $t$, predict the eye
  state label at time $t + 1$.
\end{quote}

The input is represented as a flattened feature vector of shape
$(C \times T,)$, where $C = 16$ is the number of channels and $T$ is the
window length in samples. Flattening follows row major order:

\[
  \mathbf{x} =
  \bigl[
    x_{0,0},\; x_{0,1},\; \ldots,\; x_{0,T 1},\;
    x_{1,0},\; \ldots,\;
    x_{C 1,T 1}
  \bigr]
\]

where $x_{c,\ell}$ is the value of channel~$c$ at lag~$\ell$ before the
target. This ordering gives each
feature a unique (channel, temporal lag) interpretation, which is later
exploited for feature importance analysis.

The sliding window builder constructs $N - T$ labelled samples from an
$N$-sample recording:

\begin{align*}
  \mathbf{X}_{\mathrm{win}}[i] &= \mathrm{flatten}\!\left(
      \mathbf{X}_{\mathrm{scaled}}[i:\,i+T]
    \right) \in \mathbb{R}^{C \times T} \\
  y_{\mathrm{win}}[i] &= y_{\mathrm{raw}}[i + T]
\end{align*}

\subsubsection{Model}

A \textbf{Random Forest} ensemble (\texttt{sklearn.ensemble.RandomForestClassifier})
is used as the classifier. Key motivations:

\begin{itemize}[leftmargin=1.5em]
  \item Handles high dimensional tabular inputs (up to
        $C \times T = 2\,000$ features) without architectural design.
  \item Interpretable via Gini feature importance, enabling inspection of
        which channels and temporal lags contribute most.
  \item CPU parallelisable (\texttt{n\_jobs=-1}), making online inference
        feasible on a Raspberry~Pi without a GPU (\textbf{ca veut dire, letape apres lentreinement du modele}).
  \item Produces well calibrated class probabilities via
        \texttt{predict\_proba}, suitable for continuous monitoring.
\end{itemize}

\subsubsection{Hyperparameter Optimisation}

Hyperparameter search is performed with \textbf{Optuna} (40 trials, TPE
sampler, MedianPruner). The window length~$T$ is treated as a
hyperparameter alongside the forest parameters because it directly
determines the number of input features ($C \times T$) and the temporal
context available to the model.

The full search space is shown in Table~\ref{tab:search_space}.

\begin{table}[ht]
\centering
\caption{Optuna search space for the causal Random Forest.}
\label{tab:search_space}
\begin{tabular}{lll}
  \toprule
  Hyperparameter       & Range / Choices            & Notes \\
  \midrule
  $T$                  & 4 – 128 (log scale)        & 31\,ms – 1\,s look back at 128\,Hz \\
  \texttt{n\_estimators}    & 50 – 500 (log scale)       & Number of trees \\
  \texttt{max\_depth}       & None, 5, 10, 20, 30        & None = fully grown \\
  \texttt{min\_samples\_split} & 2 – 20                  & Min samples to split a node \\
  \texttt{min\_samples\_leaf}  & 1 – 10                  & Min samples per leaf \\
  \texttt{max\_features}    & \texttt{sqrt}, \texttt{log2}, 0.1, 0.3 & Features per split \\
  \texttt{class\_weight}    & \texttt{balanced}, None    & Class imbalance handling \\
  \bottomrule
\end{tabular}
\end{table}

\textbf{Inner cross validation} uses 3 fold stratified k fold. After each
fold, the intermediate AUC is reported to the Optuna pruner; trials falling
below the median of completed trials after the first fold are pruned
(MedianPruner, 10 warm up trials). The objective function maximises the
mean ROC AUC across the three inner folds.

\subsubsection{Best Configuration}

The best trial (trial~\#21 out of 40) achieved an inner CV AUC of
\textbf{0.9996} with the configuration shown in
Table~\ref{tab:best_params}.

\begin{table}[ht]
\centering
\caption{Best hyperparameter configuration found by Optuna.}
\label{tab:best_params}
\begin{tabular}{ll}
  \toprule
  Parameter & Value \\
  \midrule
  $T$ & 125 samples ($\approx$\,977\,ms look back) \\
  \texttt{n\_estimators}    & 50 \\
  \texttt{max\_depth}       & 30 \\
  \texttt{min\_samples\_split} & 13 \\
  \texttt{min\_samples\_leaf}  & 3 \\
  \texttt{max\_features}    & \texttt{sqrt} \\
  \texttt{class\_weight}    & None \\
  \bottomrule
\end{tabular}
\end{table}

With $T = 125$ and $C = 16$, the final feature vector has
$\mathbf{2\,000}$ dimensions.

\subsubsection{Final Evaluation}

The best configuration is re evaluated with \textbf{5 fold stratified
cross validation} on the full windowed dataset (21\,193 windows).
Predictions are stored at their original indices to allow reconstruction of
the full prediction time series for visualisation.
Table~\ref{tab:results} summarises the results.

\begin{table}[ht]
\centering
\caption{5 fold CV evaluation results.}
\label{tab:results}
\begin{tabular}{llll}
  \toprule
  Metric & Closed & Open & Overall \\
  \midrule
  AUC (mean $\pm$ std) & \multicolumn{3}{c}{$\mathbf{0.9997 \pm 0.0001}$} \\
  Accuracy             & \multicolumn{3}{c}{\textbf{99.05\%}} \\
  Precision            & 1.00 & 0.98 & — \\
  Recall               & 0.99 & 0.99 & — \\
  F1 score             & 0.99 & 0.99 & 0.99 \\
  Support              & 12\,878 & 8\,315 & 21\,193 \\
  \bottomrule
\end{tabular}
\end{table}

Per fold AUC values were: 0.9997, 0.9998, 0.9996, 0.9997, 0.9996 - showing
negligible variance across folds.

\subsubsection{Feature Importance Analysis}

After evaluation, the classifier is refitted on the full windowed dataset
and Gini feature importances are extracted. The importance vector of length
$C \times T$ is reshaped into a $(C, T)$ matrix to visualise the
\textbf{temporal receptive field} — which channels and which temporal lags
the forest relies on most.

The three most informative channels (summed over all lags) are shown in
Table~\ref{tab:importance} (\textbf{dans le test case ca nest pas tres important, mais juste pour illustrer quon peut isoler les canaux plus predictifs}).

\begin{table}[ht]
\centering
\caption{Top 3 channels by summed Gini importance.}
\label{tab:importance}
\begin{tabular}{cll}
  \toprule
  Rank & Channel & Total Gini Importance \\
  \midrule
  1 & ch16 & 0.2017 \\
  2 & ch15 & 0.1745 \\
  3 & ch13 & 0.1302 \\
  \bottomrule
\end{tabular}
\end{table}

Channels 13, 15, and 16 correspond to posterior/occipital electrode
positions, consistent with the well known dominance of alpha band activity
(8–13\,Hz) in that region during eye open/closed transitions - the
expected result given that eye state is used as the proxy training label.

\subsubsection{Model Persistence}

The fitted classifier, scaler, and hyperparameter metadata are serialised to
three files:

\begin{itemize}[leftmargin=1.5em]
  \item \texttt{eeg\_rf\_model.joblib} — the complete trained forest.
  \item \texttt{eeg\_scaler.joblib} — the fitted \texttt{StandardScaler}.
  \item \texttt{eeg\_rf\_params.json} — all hyperparameters plus $T$,
        \texttt{fs}, \texttt{n\_\allowbreak{}channels}, and
        \texttt{best\_\allowbreak{}cv\_\allowbreak{}auc}.
\end{itemize}

These files are copied to the \texttt{model/} directory of PiEEG server for
deployment. So each subject will have a file containing his trained model.

% =============================================================================
\section{PiEEG Server Modifications}
% =============================================================================

The following changes were made to the PiEEG server to support online
inference and ground truth annotation.

\subsection{Button Input (GPIO)}

A physical push button connected to \textbf{GPIO~17} was added as a
ground truth label source, allowing the user to mark the target class
synchronously with the EEG stream during live recording.

\subsubsection{\texttt{hardware.py}}

The constant \texttt{BUTTON\_PIN = 17} was defined and the GPIO line is
claimed during device initialisation:

\begin{lstlisting}[style=pycode, caption={Button GPIO initialisation and read in \texttt{PiEEGHardware}.}]
BUTTON_PIN = 17

def _init_button(self):
    """Initialize button input pin (GPIO 17)."""
    self._button_fd = self._request_line(
        self._chip_fd, BUTTON_PIN,
        _GPIOHANDLE_REQUEST_INPUT, consumer=b"button"
    )

def read_button(self) -> int:
    """Read the button pin. Returns 1 when pressed, 0 otherwise."""
    buf = bytearray(_HANDLE_DATA_SIZE)
    fcntl.ioctl(self._button_fd, _GPIOHANDLE_GET_VALUES, buf)
    return buf[0]
\end{lstlisting}

The \texttt{close()} method was updated to release the button file
descriptor alongside the existing SPI and DRDY lines.

\subsubsection{\texttt{mock.py}}

A \texttt{read\_button()} stub was added to \texttt{MockHardware} so that
the mock and hardware backends remain interface compatible:

\begin{lstlisting}[style=pycode, caption={Button stub in \texttt{MockHardware}.}]
def read_button(self) -> int:
    return 0
\end{lstlisting}

\subsubsection{\texttt{acquisition.py}}

Both the hardware and mock acquisition loops were updated to call
\texttt{read\_button()} on every sample and include the result in each
broadcasted frame as a \texttt{"button"} field:

\begin{lstlisting}[style=pycode, caption={Button field added to every acquisition frame.}]
button = self._hw.read_button()
frame = {
    "t":        round(timestamp, 6),
    "n":        self._sample_count,
    "channels": sample,
    "button":   button,
}
self._loop.call_soon_threadsafe(self._enqueue, frame)
\end{lstlisting}

This extends the WebSocket frame protocol: every data frame now carries a
\texttt{"button"} integer field (0 or 1) alongside the \texttt{"channels"}
array. The value is also forwarded to the classifier to be stored as the
ground truth label for each prediction.

\subsection{Real Time Classifier Module}

A new \texttt{pieeg\_server/classifier.py} module implements
sample by sample online inference using the trained Random Forest model.

The \texttt{Classifier} class:

\begin{itemize}[leftmargin=1.5em]
  \item Subscribes to the \texttt{AcquisitionLoop} queue.
  \item Loads \texttt{eeg\_rf\_model.joblib}, \texttt{eeg\_scaler.joblib},
        and \texttt{eeg\_rf\_params.json} from the \texttt{model/} directory
        (configurable via \texttt{PIEEG\_MODEL\_DIR}).
  \item Maintains a rolling \texttt{deque} buffer of exactly $T$ raw channel
        frames.
  \item Once the buffer is full, stacks the $T$ frames into a
        $(T, C)$ array, applies the persisted scaler, flattens to
        $(1, C \times T)$, and calls \texttt{predict\_proba} to obtain
        $P(\text{eyes open})$ for the next sample.
  \item RF inference runs in a thread pool executor to avoid blocking the
        main asyncio event loop.
  \item Each prediction is appended to \texttt{self.probs} as
        \texttt{\{"prob": float, "true": button\}} and forwarded to a
        registered \texttt{\_on\_predict} async callback.
\end{itemize}

\begin{lstlisting}[style=pycode, caption={Core inference loop and prediction method in \texttt{Classifier}.}]
async def run(self):
    while True:
        frame = await self._queue.get()
        self._buffer.append(frame["channels"])
        if len(self._buffer) < self._T:
            continue
        prob = await asyncio.get_event_loop().run_in_executor(
            None, self._predict_from_buffer
        )
        self.probs.append({"prob": prob, "true": frame["button"]})
        if self._on_predict is not None:
            await self._on_predict(prob)

def _predict_from_buffer(self) -> float:
    X_raw    = np.stack(self._buffer, axis=0).astype(np.float32)
    X_scaled = self._scaler.transform(X_raw)        # (T, C)
    X_win    = X_scaled.flatten()[np.newaxis, :]    # (1, C*T)
    return float(self._clf.predict_proba(X_win)[0, 1])
\end{lstlisting}

\subsection{Prediction Broadcasting in the WebSocket Server}

\texttt{pieeg\_server/server.py} was extended to wire the classifier to the
live WebSocket broadcast loop.

\textbf{\texttt{enable\_predictor(classifier)}} registers the classifier
and installs the callback:

\begin{lstlisting}[style=pycode, caption={Registering the classifier with the WebSocket server.}]
def enable_predictor(self, classifier) -> None:
    self._classifier = classifier
    classifier._on_predict = self._on_prediction
\end{lstlisting}

\textbf{\texttt{\_on\_prediction(prob)}} is called after each inference.
It maintains a rolling history of the last 500 probabilities and broadcasts
a \texttt{"predict"} JSON message to all connected clients:

\begin{lstlisting}[style=pycode, caption={Prediction callback: history update and broadcast.}]
async def _on_prediction(self, prob: float) -> None:
    self._predict_history.append(round(prob, 4))
    if len(self._predict_history) > self._predict_history_maxlen:
        self._predict_history.pop(0)

    payload = json.dumps({
        "predict": {
            "prob":  round(prob, 4),
            "label": int(prob >= 0.5),
            "t":     time.time(),
        }
    })
    for ws in list(self._clients):
        try:
            await ws.send(payload)
        except websockets.ConnectionClosed:
            self._clients.discard(ws)
\end{lstlisting}

The \textbf{welcome message} sent to newly connected clients was extended
with \texttt{"predict\_\allowbreak{}enabled"} and
\texttt{"predict\_\allowbreak{}history"} fields, so
that a client connecting mid session immediately receives the last 250
probability values and can render the time series without waiting for new
predictions.

\subsection{Probability Time Series in the Dashboard}

The React/TypeScript dashboard was updated to display the classifier output
as a real time probability time series.

\subsubsection{\texttt{useEEG.ts} — ring buffer}

A circular \texttt{Float32Array} buffer is maintained for incoming
prediction messages:

\begin{lstlisting}[style=tscode, caption={Ring buffer for probability samples in \texttt{useEEG.ts}.}]
const predictBufRef   = useRef(new Float32Array(PRED_BUFFER_SIZE));
const predictWriteRef = useRef(0);
const predictCountRef = useRef(0);

if ("predict" in msg) {
    const p = (msg as any).predict;
    const prob = typeof p === "number" ? p : p?.prob ?? null;
    if (prob !== null) {
        const buf = predictBufRef.current;
        buf[predictWriteRef.current] = prob;
        predictWriteRef.current =
            (predictWriteRef.current + 1) % PRED_BUFFER_SIZE;
        if (predictCountRef.current < PRED_BUFFER_SIZE)
            predictCountRef.current++;
    }
}
\end{lstlisting}

Historical values received in the welcome message are pre loaded into the
same ring buffer on connect.

\subsubsection{\texttt{\_probabilityPlot.tsx} — canvas renderer}

The function \texttt{drawProbabilityPlot} reads the most recent
\texttt{plotW} samples in chronological order and draws:

\begin{itemize}[leftmargin=1.5em]
  \item Y axis ticks at 0.0, 0.5, and 1.0.
  \item A blue line trace (\texttt{\#58a6ff}) of $P(\text{eyes open})$
        over time.
  \item A filled circle at the current (rightmost) probability value.
  \item A numeric label showing the latest probability to three decimal
        places.
\end{itemize}

\subsubsection{\texttt{Spectrogram.tsx} — display mode}

The existing \texttt{Spectrogram} component was extended with a
\texttt{displayMode} prop (\texttt{"auto"} $|$ \texttt{"prob"} $|$
\texttt{"spec"}). When the mode is \texttt{"prob"} (or \texttt{"auto"}
with live predictions available), the probability renderer replaces the
spectrogram:

\begin{lstlisting}[style=tscode, caption={Display mode switch in \texttt{Spectrogram.tsx}.}]
const wantProb =
    displayMode === "prob" ||
    (displayMode === "auto" && hasPredict &&
     eegData.predictCount!.current > 1);
if (wantProb)
    drawProbabilityPlot(ctx, w, h, eegData);
\end{lstlisting}

\subsubsection{\texttt{App.tsx} — Prediction toggle button}

A \textbf{Prediction} toggle button was added to the sidebar's
\emph{Analysis} section. Clicking it forces the spectrogram panel open and
passes \texttt{displayMode="prob"} to it:

\begin{lstlisting}[style=tscode, caption={Prediction toggle in the sidebar (\texttt{App.tsx}).}]
<button
    className={`btn${showProbability ? " active" : ""}`}
    onClick={() => {
        setShowSpectrogram(true);
        setShowProbability((v) => !v);
    }}
    title="Toggle prediction time series in analysis panel"
>
    Prediction
</button>
\end{lstlisting}

% =============================================================================
\section{Experimental Results - Eye Open/Closed Test Case}
% =============================================================================

This section presents the results of applying the causal Random Forest
pipeline to the eye open/closed dataset, which serves as the validation
proxy for the button press prediction task. All figures were generated
directly from the \texttt{eeg\_causal\_rf.ipynb} notebook.

% ── Optuna search ────────────────────────────────────────────────────────────
\subsection{Hyperparameter Search}

Figure~\ref{fig:optuna} shows the outcome of the 40 trial Optuna search.
The left panel plots each completed trial as a function of the look back
window $T$ (in milliseconds) against the inner 3 fold CV AUC. Performance
is already high across a wide range of $T$ values, but saturates and
becomes most consistent for $T \gtrsim 600$\,ms, with the optimal
$T = 125$ samples ($\approx$\,977\,ms) marked by the dashed red line. The
right panel shows the optimisation history: the best so far curve (gold)
converges rapidly and stabilises after roughly 25 trials.

\begin{figure}[ht]
  \centering
  \includegraphics[width=\textwidth]{plot_optuna.png}
  \caption{Optuna hyperparameter search results. \textit{Left:} look back
    duration $T$ vs.\ inner CV AUC for each completed trial; the dashed
    line marks the selected $T = 125$ samples (977\,ms).
    \textit{Right:} optimisation history showing per trial AUC and the
    best so far curve.}
  \label{fig:optuna}
\end{figure}

% ── Confusion matrix ─────────────────────────────────────────────────────────
\subsection{Classification Performance}

Figure~\ref{fig:cm} shows the confusion matrix obtained from 5 fold
stratified cross validation (21\,193 windows). The classifier achieves
near perfect separation: only $\sim$200 windows are misclassified out of
more than 21\,000. The slight asymmetry between the two error cells
(more open predicted as closed than the reverse) is consistent with the
61\%/39\% class imbalance.

\begin{figure}[ht]
  \centering
  \includegraphics[width=0.48\textwidth]{plot_confusion_matrix.png}
  \caption{Confusion matrix from 5 fold CV ($N = 21\,193$ windows,
    $T = 125$ samples). Overall accuracy: 99.05\%.}
  \label{fig:cm}
\end{figure}

% ── Feature importance heatmap ───────────────────────────────────────────────
\subsection{Feature Importance}

Figure~\ref{fig:heatmap} shows the Gini importance reshaped into a
$16 \times 125$ heatmap (channel $\times$ temporal lag). The brightest
cells are concentrated in the posterior channels (ch13–ch16) and at
short lags (recent samples), indicating that the forest primarily relies
on recent occipital activity to make its predictions. This is consistent
with the well known alpha band desynchronisation signature of eye
opening in posterior regions.

\begin{figure}[ht]
  \centering
  \includegraphics[width=\textwidth]{plot_importance_heatmap.png}
  \caption{Feature importance heatmap. Rows are EEG channels; columns are
    temporal lags (0 = most recent sample). Colour encodes Gini importance
    — brighter regions are more discriminative. The highlighted cell marks
    the single most important (channel, lag) pair.}
  \label{fig:heatmap}
\end{figure}

Figure~\ref{fig:ch_importance} shows the total importance summed across
all lags for each channel. Channels ch16, ch15, and ch13 account for more
than 50\% of the total importance, while frontal channels (ch1–ch4)
contribute comparatively little.

\begin{figure}[ht]
  \centering
  \includegraphics[width=0.85\textwidth]{plot_channel_importance.png}
  \caption{Per channel importance (sum of Gini importance over all
    $T = 125$ lags). The most important channel (ch16) is highlighted
    in gold.}
  \label{fig:ch_importance}
\end{figure}

% ── Prediction time series ───────────────────────────────────────────────────
\subsection{Prediction Time Series}

Figure~\ref{fig:timeseries} shows a 1\,500 sample excerpt of the
prediction output alongside the 16 raw EEG channels. The two bar rows at
the top encode the true label (white) and the predicted label (coral), with
gold segments marking misclassifications. The EEG signal rows use a light
green background tint where eyes are truly open. The classifier tracks the
true label with high fidelity across the excerpt, and misclassifications
are sparse and brief - typically occurring at the onset or offset of an
eye state transition.

\begin{figure}[ht]
  \centering
  \includegraphics[width=\textwidth]{plot_signals_with_bars.png}
  \caption{1\,500 sample excerpt of the causal RF predictions overlaid on
    all 16 EEG channels. \textit{Top two rows:} true label (white bar) and
    predicted label (coral bar); gold segments indicate misclassifications.
    \textit{Lower rows:} z scored EEG channels with light green background
    where eyes are truly open.}
  \label{fig:timeseries}
\end{figure}

% 

\end{document}
